% !PRINTED 154
\weberpage{154}
Um drei Substitutionen derselben Dimension, $C$, $B$, $A$, zusammenzusetzen, muss man die Ausdrücke (10) für $z$ in eine neue lineare Substitution

\weq{12}{(u) = C(z)}

einführen und die Variablen $u$ durch die $x$ ausdrücken:

\weq{13}{(u) = CBA(x).}

Da es offenbar gleichgültig ist, ob man die Ausdrücke (10) in (12) einführt, oder ob man zuerst $z$ nach (9) durch $y$ und dann $y$ nach (4) durch $x$ ausdrückt, so gilt für diese Composition das \gesperrt{associative Gesetz}:

\weq{14}{C(BA) = (CB)A = CBA,}

was sich auch leicht durch Rechnung bestätigen lässt, wenn man nach (11) die Elemente von $C(BA)$ und $(CB)A$ bildet. Man findet für beide den Ausdruck:

\[
\sum^i \sum^h c_i^{(k)}\, b_h^{(i)}\, a_l^{(h)}.
\]

Wir sprechen also den Satz aus:

\par
3.\enspace\gesperrt{Bei der Zusammensetzung der linearen Substitutionen gilt das associative, aber nicht immer das commutative Gesetz.}

Eine Multiplication, bei der alle Multiplicatoren einander gleich sind, also

\[
N = \begin{pmatrix} \nu, & 0, & \ldots & 0 \\ 0, & \nu, & \ldots & 0 \\ \cdot & \cdot & \cdot & \cdot \\ 0, & 0, & \ldots & \nu \end{pmatrix}
\]

soll eine \gesperrt{Aehnlichkeitssubstitution} genannt werden, und zwei Substitutionen, die, wie $A$ und $AN$ oder $A$ und $NA$, durch Zusammensetzung mit einer Aehnlichkeitssubstitution aus einander abgeleitet werden können, heissen \gesperrt{ähnliche Substitutionen}. Aus dem Gesetze der Composition ergeben sich sofort die Sätze:

\par
4.\enspace\gesperrt{Eine Aehnlichkeitssubstitution ist mit jeder Substitution $A$ derselben Dimension vertauschbar; zwei Aehnlichkeitssubstitutionen, mit einander componirt, geben wieder eine Aehnlichkeitssubstitution, und zwei mit einer dritten ähnlichen Substitutionen sind auch unter einander ähnlich.}
